\subsection{A universal geometric decision criterion for zero Euler number}

\begin{statement}{C.2.UN.39}{Proved}{Curve genera, triple points and the finite monodromy constant decide zero Euler number}
Let $f:U\to\Delta$ be an actual proper holomorphic
two-torus degeneration with smooth total threefold,
and let $S=f^{-1}(0)_{\rm red}$. Choose an embedded
resolution $\pi:\widetilde U\to U$ as a sequence
of blowups with smooth centres contained in the
successive central supports. Write its full
central divisor as $\sum_jN_jE_j$ with simple
normal crossings. Let $\mathcal D$ be the set
of connected double curves of $\bigcup_jE_j$,
and let $t_3$ be its number of triple points.
Count every point centre at every stage of
$\pi$, obtaining an integer $p_\pi$, and list
every connected smooth curve centre as $C_\alpha$.
Define the integer
\[
 I_\pi=
     \sum_{D\in\mathcal D}\bigl(g(D)-1\bigr)+t_3
       +p_\pi+\sum_\alpha\bigl(1-g(C_\alpha)\bigr).
                                                               \tag{C.ZC.1}
\]
It is independent of the chosen embedded
resolution, including additional vertical
blowups used to obtain another such resolution.

Let $T$ be the actual monodromy on $H^1(F,\Z)$
of a nearby torus, and let $M$ be the order of
its semisimple part. Put
\[
             j(T)=\frac1M\sum_{k=1}^{M}\det(I-T^k).
                                                               \tag{C.ZC.2}
\]
Then
\[
                e(S)=j(T)-2I_\pi,\qquad
                e(S)=0\quad\Longleftrightarrow\quad
                         I_\pi=\frac{j(T)}2 .        \tag{C.ZC.3}
\]
The constant $j(T)$ belongs to $\{0,4,6,8\}$.
For all nineteen monodromy polynomials its
value and the required zero index are
\[
 \begin{array}{c|r|r|r}
\det(xI-T)&M&j(T)&I_\pi\text{ for }e(S)=0\\ \hline
\Phi_{1}^{4}&1&0&0\\
\Phi_{1}^{2} \Phi_{2}^{2}&2&0&0\\
\Phi_{1}^{2} \Phi_{3}&3&0&0\\
\Phi_{1}^{2} \Phi_{4}&4&0&0\\
\Phi_{1}^{2} \Phi_{6}&6&0&0\\
\Phi_{2}^{4}&2&8&4\\
\Phi_{3}^{2}&3&6&3\\
\Phi_{2}^{2} \Phi_{4}&4&4&2\\
\Phi_{4}^{2}&4&6&3\\
\Phi_{5}&5&4&2\\
\Phi_{2}^{2} \Phi_{3}&6&4&2\\
\Phi_{2}^{2} \Phi_{6}&6&4&2\\
\Phi_{3} \Phi_{6}&6&4&2\\
\Phi_{6}^{2}&6&6&3\\
\Phi_{8}&8&4&2\\
\Phi_{10}&10&4&2\\
\Phi_{3} \Phi_{4}&12&4&2\\
\Phi_{4} \Phi_{6}&12&4&2\\
\Phi_{12}&12&4&2\\
\end{array}

                                                               \tag{C.ZC.4}
\]
Thus this is a necessary and sufficient
geometric test for the zero Euler number of
every actual fibre, using only the indicated
curve genera, triple points, resolution centres
and finite monodromy constant. Arbitrary
component and centre multiplicities are
included through the full divisor and the
actual blowup maps. The statement classifies
this numerical property of actual fibres;
it makes no occurrence assertion for abstract
genus or monodromy data and does not enumerate
analytic isomorphism classes of supports.
\end{statement}

\begin{proof}
An embedded resolution of the specified kind
exists by the analytic resolution theorem used
in UN11, with all centres over the central
support and with an isomorphism over
$\Delta^*$. Smooth divisor centres need
not be included, since their blowups are
isomorphisms. The remaining connected centres
in a threefold are compact smooth curves or
points. A disconnected centre can be counted
by its disjoint connected components; their
blowups commute. In particular every genus
and count in (C.ZC.1) is attached to an actual
compact centre or curve in this resolution.

We first compute the resolved support without
any reduction of its multiplicities. Put
\[
             E_j^\circ=E_j\setminus\bigcup_{i\ne j}E_i.
\]
The proper analytic A'Campo comparison, with
its actual nearby-cycle maps as established
in (C.MS.6), says for every positive $k$ that
\[
         \det(I-T^k)=\sum_{N_j\mid k}N_j e(E_j^\circ).
                                                               \tag{C.ZC.5}
\]
Here the full integer $N_j$ is used for each
component. Choose $L$ divisible by $M$ and
by every $N_j$. The determinant on the left
depends only on the eigenvalues of $T$,
so is periodic with period $M$ even when
its unipotent part is nontrivial. Average
(C.ZC.5) over $1\le k\le L$. Exactly $L/N_j$
of these integers are divisible by $N_j$.
It follows that
\[
 \begin{split}
 j(T)
 &=\frac1L\sum_{k=1}^{L}\det(I-T^k)\\
 &=\sum_j\frac{N_j}{L}\frac{L}{N_j}e(E_j^\circ)
   =\sum_j e(E_j^\circ).
 \end{split}                                                   \tag{C.ZC.6}
\]
This calculation includes components whose
multiplicity does not divide $M$; choosing
$L$ retains them before the averaging.

Let $\widetilde S=\bigcup_jE_j$. Its
remaining strata are the double curves
with their triple points removed, and
the triple points themselves. Every
triple point lies on exactly three
double curves. Since each connected
smooth compact double curve $D$ has
Euler number $2-2g(D)$, additivity gives
\[
 \begin{split}
 e(\widetilde S)
 &=\sum_j e(E_j^\circ)
       +\sum_{D\in\mathcal D}(2-2g(D))-3t_3+t_3\\
 &=j(T)-2\left(
           \sum_{D\in\mathcal D}(g(D)-1)+t_3\right).
 \end{split}                                                   \tag{C.ZC.7}
\]
The surface Euler numbers have disappeared:
the open surface contribution is the finite
monodromy average computed in (C.ZC.6).

We now compare the original support and
the resolved one through each actual
blowup. Let $\sigma:U'\to U_0$ blow up
a connected smooth centre $Z$ contained
in a current reduced central support
$S_0$. Its reduced inverse image is
the next central support $S'$. Away
from $Z$ the map is an isomorphism;
over every point of $Z$ its fibre is
$\mathbb P^{c-1}$, where $c$ is the
complex codimension of $Z$ in $U_0$.
The entire projective fibre belongs
to $S'$, because the centre is contained
in $S_0$. Therefore the actual proper
map, stratified by $Z$ and its complement,
gives
\[
                    e(S')-e(S_0)=(c-1)e(Z).
                                                               \tag{C.ZC.8}
\]
For a point this is $2$; for a smooth
connected curve it is $2-2g(Z)$.
This computation is valid for every
multiplicity of the full fibre along
the centre. More explicitly, if the
generic order of its local equation
along the centre is $\mu$, its pullback
has the exceptional factor of order
$\mu$. That factor changes the full
divisor coefficient but not the
projective fibre in its reduced
inverse image. No value of $\mu$
was assumed in (C.ZC.8).

Iterating the actual proper maps in
$\pi$ gives
\[
 e(\widetilde S)=e(S)+2p_\pi
                        +2\sum_\alpha(1-g(C_\alpha)).
                                                               \tag{C.ZC.9}
\]
Combining (C.ZC.7) and (C.ZC.9)
proves (C.ZC.3) with exactly the
integer in (C.ZC.1). Every such
resolution is an isomorphism over
the punctured disc and therefore
has the same actual monodromy.
The formula consequently identifies
$I_\pi$ with $(j(T)-e(S))/2$ for
each choice. This proves independence
by the actual blowup comparisons;
no comparison between unrelated
fillings is required.

The nineteen possible characteristic
polynomials and their one-period
determinants were computed in UN11.
Applying the finite average (C.ZC.2)
gives the entire table (C.ZC.4).
Equivalently, $j=0$ when the polynomial
has a factor $\Phi_1$; $j=8$ for
$\Phi_2^4$; $j=6$ for
$\Phi_3^2,\Phi_4^2,\Phi_6^2$;
and $j=4$ for the ten other rows.
The exact independent calculation
in \path{code/A_zero_curve_index.py}
checks all nineteen finite averages
against the sums of the A'Campo
signature coefficients, and writes
the displayed table and
\path{records/A_zero_curve_index.csv}.
All arithmetic is integral or rational.

For a concrete check of the geometric
terms, start from a smooth torus support
with identity monodromy. A point blowup
creates one rational double curve and
no triple point. The new support has
index $-1$ and Euler number $2$.
Follow it by the genus-three quartic
blowup of UN4. There are then five
rational double curves, one double
curve of genus three, and four triple
points. Their index is
$-5+(3-1)+4=1$, giving Euler number
$-2$. Viewed instead as a resolution
of the original smooth support, these
two centres contribute
$1+(1-3)=-1$ to the centre term,
so (C.ZC.1) returns index zero for
that original support. These exact
charts verify both the sign of the
centre correction and its compatibility
with the universal balancing operations.
\end{proof}
